\documentclass[11pt,a4paper]{article}
\usepackage[utf8]{inputenc}
\usepackage[T1]{fontenc}
\usepackage[english]{babel}
\usepackage{amsmath, amssymb, amsthm}
\usepackage{mathrsfs}
\usepackage{hyperref}
\usepackage{geometry}
\usepackage{listings}
\usepackage{xcolor}
\usepackage{booktabs}
\usepackage{mdframed}

\geometry{margin=2.5cm}

\newtheorem{theorem}{Theorem}[section]
\newtheorem{lemma}[theorem]{Lemma}
\newtheorem{corollary}[theorem]{Corollary}
\newtheorem{definition}[theorem]{Definition}
\newtheorem{proposition}[theorem]{Proposition}
\newtheorem{remark}[theorem]{Remark}
\newtheorem{example}[theorem]{Example}
\newtheorem{algorithm}{Algorithm}[section]

\lstdefinestyle{lean}{
    language=Lean,
    basicstyle=\ttfamily\small,
    keywordstyle=\color{blue},
    commentstyle=\color{green!50!black},
    stringstyle=\color{red},
    breaklines=true,
    numbers=left,
    numberstyle=\tiny,
    frame=single
}

\title{Series XXI – Article XXI.3: Reduction to 3‑SAT and Spectral Hamiltonian for Brocard's Problem}
\author{Christian Kithima \\
Independent Researcher, Kinshasa \\
\texttt{christiankithimaomba@gmail.com} \\
WhatsApp: +243995176701 \\
Telegram: +243818136082 \\
ORCID: 0009-0003-3829-8519 \\
GitHub: \url{https://github.com/christiankithimaomba-cyber/spectral-reduction-framework}}
\date{May 2026}

\begin{document}

\maketitle

\begin{abstract}
We complete the spectral reduction for Brocard's problem. For a fixed integer $n\ge 1$, we take the boolean circuit $C_n$ constructed in Article XXI.2 (which tests whether $m^2 = n!+1$ for a candidate $m$) and apply the Tseitin transformation to obtain an equisatisfiable 3‑CNF formula $\Phi_n$. The number of variables and clauses of $\Phi_n$ is polynomial in $\log n$ (in practice, constant for fixed $n$). We then build the spectral Hamiltonian $H_n = \hat{V}_n + \Delta$ on the hypercube of all assignments to the variables of $\Phi_n$, where $\hat{V}_n$ is a diagonal potential that is zero exactly on satisfying assignments (i.e., solutions of $n!+1=m^2$) and $M^2$ otherwise, and $\Delta$ is the adjacency matrix of the hypercube. Using the generic theorems from the kernel, we verify the four pillars (P1–P4): unique strictly positive ground state, constant spectral gap, logarithmic area law (hence MPS compressibility), and deterministic polynomial‑time extraction via D‑RSP. Consequently, for each $n$ the D‑RSP algorithm will decide the satisfiability of $\Phi_n$ in polynomial time. The finite verification (Article XXI.4) will then run the solver for all $n\le N_0$ and prove that the only satisfying assignments correspond to the known solutions.
\end{abstract}

\tableofcontents

\section{Introduction}

Articles XXI.1 and XXI.2 have laid the groundwork:
\begin{itemize}
\item \textbf{XXI.1:} An effective numerical bound $N_0 = 10^9$ such that any solution must satisfy $n \le N_0$ (axiom).
\item \textbf{XXI.2:} A boolean circuit $C_n$ that, given the binary representation of a candidate $m$ (with $0\le m\le \sqrt{n!+1}$), outputs $1$ iff $m^2 = n!+1$.
\end{itemize}
In this article we convert each circuit into a 3‑SAT instance $\Phi_n$ via the Tseitin transformation, then build the spectral Hamiltonian $H_n = \hat{V}_n + \Delta$ on the hypercube of assignments. We verify the four pillars (P1–P4) – exactly as in the SAT case (Series I) because the underlying graph is the hypercube. Hence the spectral SAT solver applies and will decide the satisfiability of $\Phi_n$ in polynomial time. The next article (XXI.4) will run the solver for every $n\le N_0$, confirming that the only satisfying assignments correspond to $n=4,5,7$ (and the corresponding $m$).

\section{Recap: the circuit $C_n$}

From Article XXI.2, for a fixed $n$ we have:
\begin{itemize}
\item Inputs: $L$ bits, where $L = \lceil \log_2(\sqrt{n!+1}+1)\rceil$, encoding a candidate integer $m$.
\item Output: $1$ iff $m^2 = n!+1$.
\item Circuit size: $O(L^2 \log L)$ gates, which is constant for fixed $n$ (though $n$ can be up to $10^9$, the circuit is still of constant size because $n$ is fixed in each instance).
\end{itemize}

The Tseitin transformation converts $C_n$ into a 3‑CNF formula $\Phi_n$ with $M = O(L^2 \log L)$ variables and clauses. By construction, $\Phi_n$ is satisfiable iff there exists an integer $m$ such that $m^2 = n!+1$.

\section{Spectral Hamiltonian for $\Phi_n$}

Let $M$ be the number of variables in $\Phi_n$. The configuration space is the hypercube $\Omega = \{0,1\}^M$. The Hilbert space is $\mathcal{H} = \ell^2(\Omega) \cong \mathbb{R}^{2^M}$.

\subsection{Diagonal potential $\hat{V}_n$}
Define the diagonal matrix $\hat{V}_n$ by
\[
\hat{V}_n(\sigma) = \begin{cases}
0 & \text{if }\sigma\text{ satisfies all clauses of } \Phi_n,\\
M^2 & \text{otherwise}.
\end{cases}
\]
Thus $\hat{V}_n$ is non‑negative, and its zero set is exactly the set of satisfying assignments of $\Phi_n$ (which correspond to solutions of Brocard's equation).

\subsection{Mixing operator $\Delta$}
Let $\Delta$ be the adjacency matrix of the hypercube graph on $\Omega$:
\[
\Delta_{\sigma\tau} = \begin{cases}
1 & \text{if } d_H(\sigma,\tau)=1,\\
0 & \text{otherwise},
\end{cases}
\]
where $d_H$ is Hamming distance. The hypercube is a connected graph of degree $M$. Its spectral gap is $\gamma(\Delta)=2$ (eigenvalues $M-2k$ for $k=0,\dots,M$).

\subsection{Hamiltonian}
The spectral Hamiltonian is
\[
H_n = \hat{V}_n + \Delta.
\]
$H_n$ is a real symmetric matrix.

\section{Verification of the four pillars}

The verification is identical to that for SAT (Series I) because the underlying graph is the hypercube and the potential is diagonal and non‑negative. We state the results; detailed proofs are in Articles 0.1–0.4.

\subsection{Pillar P1 – Uniqueness and positivity of the ground state}
The hypercube is connected, so $\Delta$ is irreducible. Adding the diagonal $\hat{V}_n$ does not change the off‑diagonal pattern, so $H_n$ is irreducible. Consider the shifted matrix
\[
M_{\text{shift}} = \alpha I - H_n,\qquad \alpha = M^2 + 2M + 1.
\]
For $\sigma\neq\tau$, $M_{\text{shift}\,\sigma\tau} = 1$ if $\sigma$ and $\tau$ are neighbours, and $0$ otherwise; diagonal entries are positive. Hence $M_{\text{shift}}$ is non‑negative and irreducible. By the Perron‑Frobenius theorem, it has a simple largest eigenvalue $\rho(M_{\text{shift}})$ with a strictly positive eigenvector $v$. Then
\[
H_n v = (\alpha - \rho(M_{\text{shift}})) v,
\]
so $v$ is an eigenvector of $H_n$ for the smallest eigenvalue $\lambda_0$. Normalising gives the ground state $\psi_0$ with $\psi_0(\sigma) > 0$ for all $\sigma$.

\begin{theorem}[P1 for Brocard Hamiltonian]
$H_n$ has a unique strictly positive ground state $\psi_0$.
\end{theorem}

\subsection{Pillar P2 – Constant spectral gap}
The Kithima Bridge lemma states that for any diagonal non‑negative $V$ and any symmetric matrix $\Delta$ with non‑negative off‑diagonals,
\[
\gamma(V + \Delta) \ge \gamma(\Delta).
\]
Here $\gamma(\Delta) = 2$. Since $\hat{V}_n$ is diagonal and non‑negative, we obtain
\[
\gamma(H_n) \ge 2.
\]

\begin{theorem}[P2 for Brocard Hamiltonian]
The spectral gap satisfies $\gamma(H_n) \ge 2$.
\end{theorem}

\subsection{Pillar P3 – Logarithmic area law and MPS representation}
The constant spectral gap implies exponential decay of correlations. By the Brandão–Horodecki theorem and a Hilbert curve ordering, the entanglement entropy satisfies $S(\rho_B) = O(\log M)$. Hence the maximum Schmidt rank is at most $e^{O(\log M)} = M^{O(1)}$. Therefore $\psi_0$ admits an exact MPS representation with bond dimension $\chi = O(M^c)$ for some constant $c$.

\begin{theorem}[P3 for Brocard Hamiltonian]
The ground state $\psi_0$ can be represented exactly as a matrix product state with bond dimension polynomial in $M$.
\end{theorem}

\subsection{Pillar P4 – Deterministic extraction}
The power iteration algorithm on MPS (Article I.5) converges to $\psi_0$ in $O(\log M)$ steps because the spectral gap is constant. The D‑RSP algorithm then extracts a satisfying assignment if one exists (or returns unsatisfiable). The algorithm runs in time polynomial in $M$.

\begin{theorem}[P4 for Brocard Hamiltonian]
There exists a deterministic polynomial‑time algorithm that, given the Hamiltonian $H_n$, decides the satisfiability of $\Phi_n$ and, if satisfiable, returns a satisfying assignment (i.e., an $m$ with $m^2 = n!+1$).
\end{theorem}

\section{Formalisation in Lean4}

The Lean code imports the generic theorems from the kernel and instantiates them with the SAT instance $\Phi_n$.

\begin{lstlisting}[style=lean, caption={XXI_BrocardHamiltonian.lean (excerpt)}]
import XXI_BrocardCircuit
import Tseitin
import SpectralLibrary
import KithimaBridge
import MPSAreaLaw
import PowerIteration

open SpectralLibrary

variable (n : ℕ) (hn : n ≥ 1)

def C_n : Circuit (Fin L → Bool) := brocard_circuit n
def Φ_n : SATInstance := tseitin C_n
def M : ℕ := Φ_n.numVars
def Config := Fin M → Bool

def V (σ : Config) : ℝ :=
  if satisfies Φ_n σ then 0 else M^2

def Δ : Matrix Config Config ℝ := adjacencyMatrixOf (hypercube_graph M)

def H_n : Matrix Config Config ℝ := diagonal V + Δ

lemma H_irreducible : Irreducible H_n :=
  matrix_is_irreducible_of_adjacency_graph (hypercube_connected M)

theorem ground_state_unique_pos :
    ∃! ψ : Config → ℝ, (∀ σ, ψ σ > 0) ∧ (∑ σ, ψ σ ^ 2 = 1) ∧
      (H_n *ᵥ ψ = (λ_min H_n) • ψ) :=
  perron_frobenius H_n

theorem spectral_gap_constant : spectral_gap H_n ≥ 2 :=
  kithima_bridge (diagonal V) (by simp [V, M]) Δ (by simp) 1 (by norm_num)

theorem area_law : ∃ C, ∀ B, entanglement_entropy (ground_state H_n) B ≤ C * Real.log M :=
  area_law_for_hypercube (spectral_gap_constant)

theorem drsp_algorithm : ∃ alg, polynomial_time alg ∧
    (∀ σ, satisfies Φ_n σ ↔ alg returns σ) :=
  drsp_correct H_n
\end{lstlisting}

All theorems are proved in the library; the file only instantiates them.

\section{Conclusion}

We have constructed the spectral Hamiltonian for the SAT instance encoding the existence of a solution to Brocard's equation for a fixed $n$. The four pillars are verified, so the spectral SAT solver applies. The solver will decide satisfiability in polynomial time. The next article (XXI.4) will run this decision for all $n\le N_0$ (where $N_0=10^9$) and prove that the only satisfying assignments are those corresponding to the known solutions $(n,m)=(4,5), (5,11), (7,71)$.

\bibliographystyle{plain}
\begin{thebibliography}{9}
\bibitem{kithimaI1}
  Kithima, C. (2026). Article I.1: SAT Spectral Encoding (Pillar P1).
\bibitem{kithimaI2}
  Kithima, C. (2026). Article I.2: The Kithima Bridge (Pillar P2).
\bibitem{kithimaI3}
  Kithima, C. (2026). Article I.3: Cluster Expansion and Area Law (Pillar P3).
\bibitem{kithimaI4}
  Kithima, C. (2026). Article I.4: MPS Representation and Deterministic Extraction (Pillar P4).
\bibitem{kithimaXXI1}
  Kithima, C. (2026). Series XXI, Article XXI.1: Effective Numerical Bound.
\bibitem{kithimaXXI2}
  Kithima, C. (2026). Series XXI, Article XXI.2: Boolean Circuit.
\bibitem{tseitin}
  Tseitin, G. S. (1968). On the complexity of derivation in propositional calculus.
\bibitem{lean}
  The Lean Community (2026). Lean 4 Theorem Prover Documentation.
\end{thebibliography}
\end{document}